\begin{center}
\LARGE{\bfseries{Abstract}}
\end{center}
\setstretch{1.5}

The rapid advancement of Large Language Models (LLMs) has opened new frontiers in automated content synthesis; however, traditional stateless single-prompt execution models are constrained by high hallucination rates, context congestion, and fragility under system timeouts or network failures. This thesis presents the \textbf{AI Content Factory}, a stateful, multi-agent content orchestration and evaluation engine that shifts the generative paradigm from static single-turn prompt layouts to active, cooperative cognitive workflows. By modeling the content generation cycle as a stateful cyclic directed graph, the proposed system guarantees structural coherence, source grounding, and multi-channel asset integration. The core novelty of the engine lies in its ability to persist execution snapshots at every state transition, enabling reliable crash recovery and seamless human-supervised workflows for high-volume content operations.

The system is architected using a decoupled client-server model, utilizing \textbf{LangGraph} and \textbf{FastAPI} on the backend, paired with an interactive \textbf{React (Vite)} frontend dashboard. Information is gathered through web research queries (Tavily API) and semantic document parsing (Advanced RAG using cosine-similarity vector embeddings), which are filtered through an input safety node (Topic Guard). The engine features a database-backed Human-in-the-Loop (HITL) interrupt mechanism that pauses graph execution at the outline planning stage, allowing users to modify subheading parameters before launching high-resource parallel writing worker nodes. Once the sections are written and merged, a Quality Control node audits the draft against the research evidence, triggering iterative corrections via a Revision agent. The finalized text is optimized for search engines (SEO Optimizer) and dispatched to media synthesis modules, where MoviePy compiles subtitled stock videos from Pexels API queries and Gemini's native audio output API synthesizes high-fidelity solo podcasts.

To validate output quality, the system executes an automated evaluation pipeline (LLM-as-a-judge) incorporating localized G-Eval prompts and DeepEval scores across coherence, relevance, factual accuracy, and tone alignment. Testing indicates that the parallel worker dispatch reduces writing latency by over 40\% compared to sequential generation pipelines, while the database checkpointer system achieves 100\% state recovery success from forced network timeouts. The cost-benefit analysis demonstrates that generating a complete multi-modal marketing campaign (article, podcast, video, and social media copy) using the AI Content Factory costs approximately \$0.11, representing a cost reduction of over 99.9\% compared to traditional agency timelines. This project establishes a reliable, scalable, and audit-ready framework for next-generation automated media synthesis in both educational and enterprise domains.
